\documentclass[11pt,a4paper]{article}
\usepackage[T1]{fontenc}
\usepackage[utf8]{inputenc}
\usepackage{lmodern}
\usepackage[margin=27mm,headheight=14pt]{geometry}
\usepackage{amsmath,amssymb,amsthm,mathtools,mathrsfs}
\usepackage{microtype,booktabs,longtable,array}
\usepackage{xcolor,xurl}
\definecolor{ink}{RGB}{28,56,80}
\usepackage[colorlinks=true,allcolors=ink]{hyperref}
\usepackage{fancyhdr}
\pagestyle{fancy}\fancyhf{}
\fancyhead[L]{\small Polylogarithmic Lambert inversion}
\fancyhead[R]{\small ProveIt research contribution}
\fancyfoot[C]{\thepage}
\setlength{\emergencystretch}{2em}
\numberwithin{equation}{section}
\allowdisplaybreaks[2]
\newtheorem{theorem}{Theorem}[section]
\newtheorem{proposition}[theorem]{Proposition}
\newtheorem{lemma}[theorem]{Lemma}
\newtheorem{corollary}[theorem]{Corollary}
\theoremstyle{definition}\newtheorem{definition}[theorem]{Definition}
\theoremstyle{remark}\newtheorem{remark}[theorem]{Remark}
\newcommand{\Li}{\operatorname{Li}}
\newcommand{\Res}{\operatorname{Res}}
\newcommand{\Disc}{\operatorname{Disc}}
\newcommand{\PV}{\operatorname{PV}}
\newcommand{\arsinh}{\operatorname{arsinh}}
\newcommand{\artanh}{\operatorname{artanh}}
\newcommand{\Cat}{\operatorname{Cat}}
\newcommand{\dd}{\,\mathrm d}
\newcommand{\eps}{\varepsilon}
\newcommand{\ii}{\mathrm i}
\newcommand{\C}{\mathbb C}\newcommand{\R}{\mathbb R}
\newcommand{\Q}{\mathbb Q}\newcommand{\N}{\mathbb N}
\newcommand{\F}{\mathcal F}\newcommand{\CC}{\mathcal C}
\newcommand{\HH}{\mathcal H}
\newcommand{\B}{\widehat{\mathcal R}}
\newcommand{\Sc}{\mathcal S}
\newcommand{\ua}{\mathfrak a}\newcommand{\uz}{\mathfrak z}
\hypersetup{pdftitle={Double Resonance and Nested Stokes Geometry in Polylogarithmic Lambert Inversion},pdfauthor={Research contribution prepared for Vladimir Reshetnikov with OpenAI assistance}}
\begin{document}
\begin{center}
{\small\scshape Analysis / Polylogarithms / Transseries}\par\vspace{10pt}
{\LARGE\bfseries Double Resonance and Nested Stokes Geometry}\par\vspace{5pt}
{\Large\bfseries in Polylogarithmic Lambert Inversion}\par\vspace{12pt}
{\normalsize Prepared for Vladimir Reshetnikov with OpenAI assistance}\par
{\small 12 October 2026}
\end{center}
\begin{abstract}
We study the symmetric specialization
$\F_\eps(t)=\sum_{a,b\ge1}(a/b)^\eps e^{-abt}/(ab)$
of the established two-index Lambert family. Three Mellin poles coalesce
at $\eps=0$, producing a quadratic logarithmic core. A centered Gamma--zeta
normalization gives every resonant order jet and removes its next-to-leading
logarithmic degree. The trivial zeros suppress every odd perturbative power
beyond the linear term; the second order jet restores that entire parity
sector, with an explicit arithmetic Borel kernel and a Stokes jump exactly
$\pi^2/2$ times the unperturbed jump.

The principal inverse contribution is a uniform arcsinh crossover when
$\eps=\lambda/X$ and the target is $X^2/2$. We compute its first three
corrections and the displacement of the two complex folds of its singular
core. At exact resonance we give a finite formula for every inverse-sector
coefficient, prove alternating coefficientwise signs, determine exact
valuations and Catalan highest-degree terms, and prove zero ordinary
radius but exact Borel radius $4\pi^2$ for every positive fixed slow
parameter. An elementary Mellin proof fixes the exact real completion.
Lagrange inversion then transports the entire divisor-weighted Stokes jump
to the inverse. In the physical variable the jump is doubly exponentially
small. The arithmetic average of the lateral inverses differs from the
real inverse by a second-action term whose leading coefficient simplifies
to $-1$. Every asserted theorem has a conventional proof; finite symbolic
checks, numerical diagnostics, a scoped source audit, and further research
questions accompany the article.
\end{abstract}
\noindent\textbf{Status.} Unrefereed mathematical research with explicit proofs;
not a Lean formalization. The Lambert family and its general resurgent
infrastructure are prior work. The explicit resonant inverse and jet
refinements are proposed contributions, not certified claims of worldwide
priority. No repository files were modified.
\clearpage
\tableofcontents
\clearpage

\section{Research setting, antecedents, and principal results}
\label{sec:intro}

\subsection{Why this particular intersection}
A single polylogarithm near its unit argument has a locally convergent
power--log expansion after the appropriate pole cancellation. Summing
polylogarithms on a multiplicative lattice changes the analytic class.
The resulting Lambert-type function has a factorially divergent
perturbative tail and an arithmetic lattice of Borel singularities.
Nonlinear inversion then produces two distinct exponential levels:
ordinary sectors of size $e^{-nX}$ and Stokes sectors of size
$\exp(-4\pi^2 e^{X-\gamma})$.

The distinction is consequential for the two ProveIt projects. The
Polylogarithms sources develop normalized Gamma/zeta order bridges; the
Transseries sources distinguish exact cores, formal reversion, analytic
charts, and residual transport. This article connects those operations
without importing a convergence assertion from the wrong analytic class.
The inspected bridge chapter explicitly regularizes trivial zeros before
taking logarithms, and the transseries overview explicitly distinguishes
coefficient-form from operator-form Lagrange inversion
\cite{repo-bridge,repo-transseries}.

\subsection{Prior work and the claimed scope of this contribution}
The family considered here is not new. With trivial characters and
$q=e^{-t}$ it is exactly
\[
 \F_\eps(t)=\Xi_{1-\eps,1+\eps}(1,1;q)
\]
in the notation of Broadhurst and Dorigoni \cite{BD2025}, who treat
arbitrary pairs of indices and Dirichlet characters and their full
resurgent expansions. Their 2026 overview states this general family and
its transformed nonperturbative terms explicitly \cite{BD2026}.
Dorigoni and Kleinschmidt had already developed resurgent expansions of
Lambert series and iterated Eisenstein integrals \cite{DK2021}.
Accordingly, neither the family, the Mellin method, the existence of
arithmetic Borel singularities, nor the general phenomenon of hidden
exponential contributions is claimed here as a discovery.

Our narrower proposed refinements are a centered triple-collision jet
calculus, a closed parity-jet Stokes identity, an arcsinh inverse crossover
with explicit Gamma--Stieltjes corrections, a finite rational inverse
coefficient formula with sharp structural consequences, and explicit
nonlinear transport to a nested exponential scale. The formulas are
proved directly, making their normalization and domains independently
checkable. The exact scalar median identity is treated as a useful
specialization of the established Lambert/Voronoi structure, not as a new
general transformation formula.

The library counterparts of three recent research drafts were also
compared at the level of their abstracts and selected relevant
statements: the Gamma-normalized resonant Lerch jets, the uniform
single-polylogarithm fold analysis, and the Fermi--Dirac/dilogarithm
reversion package \cite{prior-jets,prior-folds,prior-fermi}.
Their analytic tails and singular cores differ from the double Lambert
sum here. We do not claim to have extracted or audited every member of
the five incoming ZIP archives. Section~\ref{sec:audit} specifies the
boundary of the inspection.

\subsection{A preview of the exact identities}
Write $d(n)$ for the number of positive divisors of $n$, and put
\begin{equation}\label{eq:A}
 A=4\pi^2,\qquad
 C_*={\zeta(2)\over2}-\gamma^2-2\gamma_1.
\end{equation}
The real resonant function is
\[
 \F_0(t)=\sum_{n\ge1}{d(n)\over n}e^{-nt}
        =-\sum_{a\ge1}{\log(1-e^{-at})\over a},\qquad t>0.
\]
Its exact Borel completion is
\begin{equation}\label{eq:preview-median}
 \F_0(t)=\frac12(\log(1/t)+\gamma)^2+C_*-\frac t4
 -2\sum_{n\ge1}d(n)\int_0^\infty
 e^{-\xi/t}\log\left|1-\frac{\xi^2}{A^2n^2}\right|\frac{\dd\xi}{\xi}.
\end{equation}
The two lateral values of the logarithm give a jump
\begin{equation}\label{eq:preview-jump}
 J(t)=4\pi\ii\sum_{n\ge1}d(n)E_1(An/t),\qquad
 E_1(x)=\int_x^\infty{e^{-v}\over v}\dd v.
\end{equation}
If $g_0$ is the real inverse of $\F_0$, and $g_\pm$ are the nearby
inverses of $\F_0\pm J/2$, then, with $t=g_0(Y)\to0^+$,
\begin{align}
 g_+(Y)-g_-(Y)
 &\sim {\ii\over\pi}\,{t^2\over\log(1/t)+\gamma}e^{-A/t},
 \label{eq:preview-inverse-jump}\\
 \frac{g_+(Y)+g_-(Y)}2-g_0(Y)
 &\sim -{t^2\over(\log(1/t)+\gamma)^2}e^{-2A/t}.
 \label{eq:preview-bias}
\end{align}
The coefficient $-1$ in the second line is an exact simplification, not a
fitted numerical constant. Arithmetic averaging in that line is not being
identified with the general nonlinear notion of median resummation.

\section{The symmetric family and its Mellin representation}
\label{sec:mellin}

\begin{definition}
For $\Re t>0$ and $\eps\in\C$, define
\begin{align}
 D_\eps(n)&=\sum_{ab=n}(a/b)^\eps,\label{eq:divisor-weight}\\
 \F_\eps(t)&=\sum_{a,b\ge1}{(a/b)^\eps\over ab}e^{-abt}
 =\sum_{n\ge1}{D_\eps(n)\over n}e^{-nt}.
 \label{eq:family}
\end{align}
The powers use real logarithms of positive integers. All logarithms of
$t$ below use the branch real on the positive axis, unless a different
local continuation is explicitly specified.
\end{definition}

\begin{proposition}[Analyticity, symmetry, and the polylogarithm bridge]
\label{prop:family}
The family is jointly holomorphic on $\{\Re t>0\}\times\C$, even in
$\eps$, and satisfies
\begin{equation}\label{eq:polylog-bridge}
 \F_\eps(t)=\sum_{a\ge1}a^{\eps-1}\Li_{1+\eps}(e^{-at}).
\end{equation}
For real $\eps$ and $t>0$ it is positive and strictly decreasing in $t$,
with range $(0,\infty)$. For $c>|\Re\eps|$ it has the Mellin formula
\begin{equation}\label{eq:mellin}
 \F_\eps(t)=\frac1{2\pi\ii}\int_{c-\ii\infty}^{c+\ii\infty}
 \Gamma(s)\zeta(1+s-\eps)\zeta(1+s+\eps)t^{-s}\dd s.
\end{equation}
\end{proposition}
\begin{proof}
On a compact parameter set with $|\Re\eps|\le M$,
$|D_\eps(n)|\le d(n)n^M$. Exponential decay in $n$ gives normal
convergence, also after any fixed number of parameter or $t$ derivatives.
Interchanging $a$ and $b$ proves evenness. Expanding the defining
polylogarithm series proves \eqref{eq:polylog-bridge}; that series
normalization is classical \cite{DLMFpolylog}.
For real $\eps$, the summands are positive and their $t$ derivatives are
negative. The limit at infinity is zero. Pairing $(a,b)$ with $(b,a)$
gives $D_\eps(n)\ge d(n)$, while the partial sum of $1/(ab)$ over
$a,b\le M$ tends to infinity as $M\to\infty$. Thus the limit at zero is
infinite. Finally insert Mellin inversion for $e^{-abt}$. Absolute
convergence on the line follows from $c>|\Re\eps|$ and the two
absolutely convergent Dirichlet series. The exponential decay of
$\Gamma(c+\ii v)$ completes the interchange and also permits
$|\arg t|<\pi/2$.
\end{proof}

The Dirichlet series of the finite arithmetic weights is particularly
simple:
\begin{equation}\label{eq:weight-dirichlet}
 \sum_{n\ge1}{D_\eps(n)\over n^z}
 =\zeta(z-\eps)\zeta(z+\eps),\qquad \Re z>1+|\Re\eps|.
\end{equation}
For later use set
\begin{equation}\label{eq:divisor-moments}
 M_r(n)=\sum_{ab=n}\log^r(a/b).
\end{equation}
Then $M_{2r+1}(n)=0$ and
\begin{equation}\label{eq:moment-zeta}
 \sum_{n\ge1}{M_{2r}(n)\over n^z}
 =\sum_{j=0}^{2r}(-1)^j\binom{2r}{j}
 \zeta^{(j)}(z)\zeta^{(2r-j)}(z),\qquad \Re z>1.
\end{equation}
This exact identity expresses the order jets both in finite divisor
moments and in ordinary zeta derivatives.

\section{Three coalescing poles and all centered order jets}
\label{sec:collision}
Use the Stieltjes convention
\begin{equation}\label{eq:stieltjes}
 \zeta(1+z)=\frac1z+\sum_{r\ge0}{(-1)^r\gamma_r\over r!}z^r,
 \qquad \gamma_0=\gamma.
\end{equation}
For $0<|\eps|<1$ define the singular core
\begin{equation}\label{eq:core}
 \CC_\eps(L)=\Gamma(\eps)\zeta(1+2\eps)e^{\eps L}
 +\Gamma(-\eps)\zeta(1-2\eps)e^{-\eps L}
 +\zeta(1-\eps)\zeta(1+\eps).
\end{equation}
Each of the first two terms has a double pole in $\eps$; the sum does not.

\begin{theorem}[Uniform pole cancellation and complete power coefficients]
\label{thm:collision}
The apparent singularity of $\CC_\eps(L)$ at $\eps=0$ is removable,
and the result is entire in $L$ and holomorphic and even in $|\eps|<1$.
Its resonant value is
\begin{equation}\label{eq:core0}
 \CC_0(L)=\frac12(L+\gamma)^2+C_*.
\end{equation}
For every $N\ge0$, $0<\rho<1$, and $\eta>0$,
\begin{equation}\label{eq:asymptotic}
 \F_\eps(t)=\CC_\eps(\log(1/t))+
 \sum_{k=1}^{N}r_k(\eps)t^k+O(t^{N+1}),
 \qquad
 r_k(\eps)={(-1)^k\over k!}\zeta(1+\eps-k)\zeta(1-\eps-k),
\end{equation}
uniformly for $|\eps|\le\rho$ and
$|\arg t|\le\pi/2-\eta$ as $t\to0$. On smaller parameter disks the
same assertion can be differentiated any fixed number of times in
$\eps$.
\end{theorem}
\begin{proof}
The Mellin integrand has the three moving poles $s=\eps,-\eps,0$.
Their residues give \eqref{eq:core}. A fixed circle $|s|=r$ with
$\rho<r<1$ represents their sum by a contour integral. Its integrand is
holomorphic in $\eps$ on a slightly larger disk and entire in $L$.
This proves removability, uniformity, and symmetry without subtracting
large singular quantities numerically.

At $\eps=0$, multiply
\[
 \Gamma(s)=s^{-1}-\gamma+
 \tfrac12(\gamma^2+\zeta(2))s+O(s^2),\qquad
 \zeta(1+s)^2=s^{-2}+2\gamma s^{-1}+\gamma^2-2\gamma_1+O(s).
\]
The coefficients of $s^{-3},s^{-2},s^{-1}$ in the product are
$1,\gamma,\zeta(2)/2-\gamma^2/2-2\gamma_1$.
Taking the residue after multiplication by $e^{sL}$ gives
\eqref{eq:core0}.

Shift the Mellin line to $\Re s=-N-1-\delta$, with
$0<\delta<1$. The additional poles crossed are the Gamma poles
$s=-1,\ldots,-N-1$, whose residues are exactly $r_k(\eps)t^k$.
Polynomial bounds for zeta in fixed vertical strips, together with
Stirling's exponential decay and the prescribed $t$ sector, give a
uniform $O(t^{N+1+\delta})$ remaining integral. Absorbing the
$k=N+1$ term proves the displayed remainder. Choose a larger compact
parameter disk still inside $|\eps|<1$ and apply Cauchy's formula to
obtain the derivative assertions. The Gamma expansion and zeta
functional identities used here have their standard normalizations
\cite{DLMFgamma,DLMFzeta}.
\end{proof}

\subsection{A centered finite formula for every logarithmic jet}
Define the removable analytic functions
\begin{equation}\label{eq:normalizers}
 \ua(z)=z^2\Gamma(z)\zeta(1+2z)e^{-\gamma z}
       =\sum_{j\ge0}a_jz^j,\qquad
 \uz(z)=z^2\zeta(1-z)\zeta(1+z)=\sum_{j\ge0}z_jz^j.
\end{equation}
Their initial coefficients are
\begin{align}
 a_0&=\tfrac12,&a_1&=0,&a_2&=b:=\frac{\zeta(2)}4-\gamma^2-2\gamma_1,
 \label{eq:a012}\\
 a_3&=\frac43\gamma^3+4\gamma\gamma_1+2\gamma_2-\frac{\zeta(3)}6,
 &z_0&=-1,&z_2&=\delta_0:=\gamma^2+2\gamma_1.
 \label{eq:a3z2}
\end{align}
All odd $z_j$ vanish. The notation $z_j$ is a coefficient label, not a
zeta value.

\begin{theorem}[Centered triple-resonance jets]
\label{thm:centered-jets}
Let $W=L+\gamma$. For every $r\ge0$,
\begin{equation}\label{eq:centered-jets}
 \left.\partial_\eps^{2r}\CC_\eps(L)\right|_{\eps=0}
 =(2r)!\left\{2\sum_{j=0}^{2r+2}
       {a_jW^{2r+2-j}\over(2r+2-j)!}+z_{2r+2}\right\}.
\end{equation}
Its degree is exactly $2r+2$, its leading coefficient is
$1/((2r+1)(2r+2))$, and the coefficient of $W^{2r+1}$ is zero.
Every odd order jet vanishes. Formula \eqref{eq:normalizers}, together
with the convergent log-Gamma series, is a finite coefficient algorithm
in Stieltjes constants and zeta values for each fixed jet.
\end{theorem}
\begin{proof}
The centered identity is
\begin{equation}\label{eq:centered-core}
 \CC_\eps(L)=\frac{
 \ua(\eps)e^{\eps W}+\ua(-\eps)e^{-\eps W}+\uz(\eps)}{\eps^2}.
\end{equation}
The numerator is even and its constant term is $2a_0+z_0=0$.
Extracting $\eps^{2r+2}$ proves \eqref{eq:centered-jets};
$a_0=1/2$ and $a_1=0$ give the two highest-degree assertions.
For completeness,
\[
 \ua(z)=\exp\left(-2\gamma z+
       \sum_{m\ge2}\frac{(-1)^m\zeta(m)}m z^m\right)
       \left(\frac12+\gamma z-2\gamma_1z^2+2\gamma_2z^3+\cdots\right).
\]
Multiplication gives \eqref{eq:a012}--\eqref{eq:a3z2} and computes
any prescribed finite order using only finite Taylor data.
\end{proof}

Thus the centered second jet has the form
\[
 \left.\partial_\eps^2\CC_\eps(L)\right|_0
 =\frac{W^4}{12}+2bW^2+4a_3W+4a_4+2z_4.
\]
The missing cubic term is a structural consequence of the centering,
not a numerical cancellation. Combining Theorems~\ref{thm:collision}
and~\ref{thm:centered-jets} gives a complete asymptotic expansion of
$\sum_n M_{2r}(n)e^{-nt}/n$ at every fixed order $2r$.

\section{Arithmetic Borel kernels and second-jet parity restoration}
\label{sec:borel}

\subsection{Reflection before differentiation}
For $k\ge2$ the zeta functional equation gives
\begin{equation}\label{eq:reflected-r}
 r_k(\eps)=\frac{2\{1+(-1)^k\cos(\pi\eps)\}}{A^k}
 \frac{\Gamma(k-\eps)\Gamma(k+\eps)}{k!}
 \zeta(k-\eps)\zeta(k+\eps).
\end{equation}
Indeed, apply
$\zeta(1-z)=2(2\pi)^{-z}\cos(\pi z/2)\Gamma(z)\zeta(z)$
to $z=k-\eps,k+\eps$, then use the cosine product formula.
The restriction $k\ge2$, $|\eps|<1$ avoids the Dirichlet-series
boundary at $1$. The linear coefficient is retained separately:
\begin{equation}\label{eq:r1}
 r_1(\eps)=-\zeta(\eps)\zeta(-\eps).
\end{equation}

\begin{definition}[Shifted Borel normalization]
For the formal tail $\mathcal R_\eps(t)=\sum_{k\ge2}r_k(\eps)t^k$,
put
\begin{equation}\label{eq:borel-convention}
 \B_\eps(\xi)=\sum_{k\ge2}{r_k(\eps)\over\Gamma(k)}\xi^{k-1}.
\end{equation}
Thus Laplace integration has no extra factor $1/t$:
$\int_0^\infty e^{-\xi/t}\xi^{k-1}\dd\xi=\Gamma(k)t^k$.
\end{definition}

\begin{theorem}[Normally convergent hypergeometric divisor kernel]
\label{thm:hypergeometric-borel}
Let
\begin{equation}\label{eq:H}
 H_\eps(z)=\Gamma(1-\eps)\Gamma(1+\eps)
 \left\{z\,{}_2F_1(1-\eps,1+\eps;2;z)-z\right\}.
\end{equation}
For $|\eps|<1$, the Borel germ has the representation
\begin{equation}\label{eq:general-borel}
 \B_\eps(\xi)=\frac2\xi\sum_{n\ge1}D_\eps(n)
 \left\{H_\eps\left(\frac\xi{An}\right)
       +\cos(\pi\eps)H_\eps\left(-\frac\xi{An}\right)\right\}.
\end{equation}
It converges normally, jointly in $\eps$ and $\xi$, on compact
subsets of $|\eps|<1$ and
$\C\setminus(( -\infty,-A]\cup[A,\infty))$, with the value at
$\xi=0$ filled in. Parameter differentiation is legitimate there.
At resonance it reduces to the elementary logarithmic kernel
\begin{equation}\label{eq:borel0}
 \B_0(\xi)=-\frac2\xi\sum_{n\ge1}d(n)
             \log\left(1-\frac{\xi^2}{A^2n^2}\right).
\end{equation}
The Taylor radius of this germ is exactly $A$.
\end{theorem}
\begin{proof}
The coefficient of $z^k$, $k\ge2$, in $H_\eps$ is
\[
 \frac{\Gamma(k-\eps)\Gamma(k+\eps)}{k!\,\Gamma(k)}.
\]
Insert this coefficient into \eqref{eq:general-borel}, apply
\eqref{eq:weight-dirichlet}, and compare with
\eqref{eq:reflected-r}. This proves the germ identity. For large $n$,
$H_\eps(\xi/(An))=O_K(n^{-2})$ uniformly on any stated compact
set; also $|D_\eps(n)|\le d(n)n^\rho$ for some $\rho<1$.
The majorant $\sum d(n)n^{-2+\rho}$ converges. Finitely many
remaining hypergeometric terms have their usual analytic branches away
from their cuts. This proves normal convergence and the derivative claim.

At $\eps=0$,
$z\,{}_2F_1(1,1;2;z)=-\log(1-z)$, giving
$H_0(z)=-\log(1-z)-z$ and hence \eqref{eq:borel0}.
The $n=1$ term has a nonzero logarithmic singularity at $\xi=A$;
all $n\ge2$ terms are analytic nearby. The same holds at $-A$.
The germ is analytic in $|\xi|<A$, proving the exact radius.
\end{proof}

\subsection{An entire parity sector restored by the second jet}
At $\eps=0$, formula \eqref{eq:reflected-r} gives
\begin{equation}\label{eq:even-coefficients}
 r_1(0)=-\frac14,\qquad
 r_{2m}(0)=\frac{4\Gamma(2m)}{2m\,A^{2m}}\zeta(2m)^2>0,
 \qquad r_{2m+1}(0)=0\quad(m\ge1).
\end{equation}
In particular the nonzero even coefficients grow factorially. The
vanishing odd coefficients have a double, rather than simple, zero as
functions of the order parameter.

\begin{theorem}[Exact parity activation and its Stokes multiplier]
\label{thm:parity}
For every odd $k\ge3$,
\begin{equation}\label{eq:second-jet-coef}
 r_k''(0)=\frac{2\zeta'(1-k)^2}{k!}
 =\frac{2\pi^2}{A^k}\frac{\Gamma(k)^2}{k!}\zeta(k)^2>0.
\end{equation}
The Borel transform of the restored odd-power part is
\begin{equation}\label{eq:odd-jet-borel}
 \widehat{\mathcal J}_{\rm odd}(\xi)
 :=\sum_{\substack{k\ge3\\k\ {\rm odd}}}
       {r_k''(0)\over\Gamma(k)}\xi^{k-1}
 =\frac{2\pi^2}{\xi}\sum_{n\ge1}d(n)
 \left\{\artanh\left(\frac\xi{An}\right)-\frac\xi{An}\right\}.
\end{equation}
For the lateral boundary values continued from zero,
\begin{align}
 \Disc\B_0(\xi)
 &=\frac{4\pi\ii}{\xi}\sum_{An<\xi}d(n),\qquad \xi>0,
 \label{eq:disc0}\\
 \Disc\widehat{\mathcal J}_{\rm odd}(\xi)
 &=\frac{2\pi^3\ii}{\xi}\sum_{An<\xi}d(n)
 =\frac{\pi^2}{2}\Disc\B_0(\xi).
 \label{eq:discjet}
\end{align}
Consequently the corresponding lateral Laplace jumps satisfy the exact
identity
\begin{equation}\label{eq:jet-jump}
 (\Sc_+-\Sc_-)\mathcal J_{\rm odd}(t)
 =2\pi^3\ii\sum_{n\ge1}d(n)E_1(An/t)
 =\frac{\pi^2}{2}(\Sc_+-\Sc_-)\mathcal R_0(t),\qquad t>0.
\end{equation}
\end{theorem}
\begin{proof}
For odd $k\ge3$, the factor in braces in \eqref{eq:reflected-r}
is $1-\cos(\pi\eps)$, whose second derivative at zero is $\pi^2$.
The other factors are nonzero and analytic there. This proves the
second expression in \eqref{eq:second-jet-coef}. Alternatively,
$\zeta(1-k+\eps)\zeta(1-k-\eps)
=-\zeta'(1-k)^2\eps^2+O(\eps^4)$, and $(-1)^k=-1$;
this proves the first expression, including its sign.

After division by $\Gamma(k)$ the coefficient is
$2\pi^2\zeta(k)^2/(kA^k)$. The identity
$\sum_{k\ge3,\ k\,\mathrm{odd}}z^k/k=\artanh z-z$,
followed by $\zeta(k)^2=\sum d(n)n^{-k}$, proves
\eqref{eq:odd-jet-borel}. All interchanges first take place in
$|\xi|<A$; the same large-$n$ estimates as before extend them to
compact slitted domains.

We fix the signs explicitly. For $\xi>An$, approaching from the upper
half-plane makes $1-\xi^2/(A^2n^2)$ approach the negative axis from
below. Thus its logarithm has upper-minus-lower jump $-2\pi\ii$.
The minus sign in \eqref{eq:borel0} gives \eqref{eq:disc0}.
Similarly
$\artanh z=(\log(1+z)-\log(1-z))/2$ has jump $\pi\ii$ for
$z>1$. This gives \eqref{eq:discjet}. Integrating each step function
against $e^{-\xi/t}/\xi$ produces $E_1(An/t)$.
The divisor sum converges exponentially, justifying the integration.
\end{proof}

\begin{remark}
``Odd-power part'' refers to the original powers $t^k$, not to parity
in the Borel coordinate: the shifted transform uses $\xi^{k-1}$.
The theorem concerns this projected part of the second jet. It does not
assert that the \emph{whole} second order jet is $\pi^2/2$ times the
original function.
\end{remark}

\subsection{Finite all-order harmonic and zeta-derivative coordinates}
The regular Gamma product in \eqref{eq:reflected-r} satisfies
\begin{equation}\label{eq:gamma-product-jets}
 \frac{\Gamma(k-\eps)\Gamma(k+\eps)}{\Gamma(k)^2}
 =\exp\left(\sum_{j\ge1}
 \frac{\zeta(2j)-H_{k-1}^{(2j)}}j\eps^{2j}\right),
 \qquad H_m^{(r)}=\sum_{\ell=1}^m\ell^{-r}.
\end{equation}
This follows by adding the Taylor series of the two log-Gamma factors;
$\psi^{(2j-1)}(k)=(2j-1)!(\zeta(2j)-H_{k-1}^{(2j)})$.
Consequently, for every $r\ge0$ and $k\ge2$,
\begin{align}
 r_k^{(2r)}(0)
 &=\frac{2\Gamma(k)^2(2r)!}{k!A^k}[\eps^{2r}]
 \Bigg(\{1+(-1)^k\cos(\pi\eps)\}
 \exp\left(\sum_{j=1}^{r}
 {\zeta(2j)-H_{k-1}^{(2j)}\over j}\eps^{2j}\right)
 \nonumber\\[-2pt]
 &\hspace{42mm}\times
 \sum_{\ell=0}^{r}\frac{\eps^{2\ell}}{(2\ell)!}
 \sum_{j=0}^{2\ell}(-1)^j\binom{2\ell}{j}
 \zeta^{(j)}(k)\zeta^{(2\ell-j)}(k)\Bigg).
 \label{eq:all-coefficient-jets}
\end{align}
Only finitely many coefficients are used at each prescribed order.
Together with \eqref{eq:centered-jets}, this supplies both the logarithmic
core and every power coefficient of every symmetric order derivative.
No differentiation of a singular logarithm at a trivial zero is needed.

\section{An exact real completion and its arithmetic jump}
\label{sec:median}
The existence of a real function with a prescribed asymptotic expansion
does not, by itself, fix any beyond-all-orders real term. We therefore
prove the equality to the specified Lambert function, rather than infer
it from reality or conjugacy of lateral values.

\begin{theorem}[Exact logarithmic Borel completion]
\label{thm:median}
For every $t>0$, define
\[
 \Sc_\pm\mathcal R_0(t)=\int_0^\infty e^{-\xi/t}\B_0(\xi\pm\ii0)\dd\xi.
\]
The integrals exist, with integrable logarithmic boundary singularities,
and
\begin{align}
 \F_0(t)&=\CC_0(\log(1/t))-\frac t4
 +\frac12\{\Sc_+\mathcal R_0(t)+\Sc_-\mathcal R_0(t)\},
 \label{eq:median-exact}\\
 \Sc_+\mathcal R_0(t)-\Sc_-\mathcal R_0(t)
 &=J(t)=4\pi\ii\sum_{n\ge1}d(n)E_1(An/t).
 \label{eq:stokes-exact}
\end{align}
Thus the forward lateral functions are exactly
\begin{equation}\label{eq:forward-laterals}
 \F_\pm(t):=\CC_0(\log(1/t))-t/4+\Sc_\pm\mathcal R_0(t)
 =\F_0(t)\pm J(t)/2.
\end{equation}
\end{theorem}

\begin{proof}
First note an integrability estimate. For fixed $t>0$ and $a\ge1$,
\[
 \int_0^\infty e^{-\xi/t}
 \left|\log\left|1-\frac{\xi^2}{a^2}\right|\right|
 \frac{\dd\xi}{\xi}\le\frac{K_t}{a^2}.
\]
On $\xi\le a/2$ use $|\log(1-v)|\le C v$ for $0\le v\le1/4$.
On $\xi>a/2$ substitute $\xi=au$ and retain half the exponential
for integrability; the remaining factor is exponentially small in $a$.
The logarithmic singularity at $u=1$ is integrable. The imaginary step
part has an even smaller exponential integral. Summation over $a=An$
is therefore justified by $\sum d(n)n^{-2}<\infty$.

It remains to identify the real part. The elementary Mellin identity
\begin{equation}\label{eq:log-mellin}
 \int_0^\infty u^{-w-1}\log|1-u^2|\dd u
 =\frac\pi w\cot\frac{\pi w}{2},\qquad 0<\Re w<2,
\end{equation}
follows by integration by parts, interpreted symmetrically at $u=1$,
and the substitution $v=u^2$. In detail, the derivative integral is
$-w^{-1}\PV\int_0^\infty v^{-w/2}/(1-v)\dd v$;
$\PV\int_0^\infty v^{a-1}/(1-v)\dd v=\pi\cot(\pi a)$
gives \eqref{eq:log-mellin}. The latter identity also follows by
splitting at $1$ and using the digamma reflection formula.

For
\[
 K(x)=-2\int_0^\infty e^{-u/x}\log|1-u^2|\frac{\dd u}{u},
\]
Mellin inversion of the exponential and \eqref{eq:log-mellin} imply
\begin{equation}\label{eq:K-mellin}
 K(x)=\frac1{2\pi\ii}\int_{b-\ii\infty}^{b+\ii\infty}
 -\frac{2\pi}{w}\cot\frac{\pi w}{2}\Gamma(w)x^w\dd w,
 \qquad 0<b<2.
\end{equation}
Absolute integration in the intermediate Mellin variable follows from
$0<b<2$ and Gamma decay, so this step does not rely on a formal exchange.
Choose $1<b<2$, sum $d(n)K(t/(An))$, and apply
$\sum d(n)n^{-w}=\zeta(w)^2$. The real average in
\eqref{eq:median-exact} is therefore
\begin{equation}\label{eq:median-MB}
 \frac1{2\pi\ii}\int_{b-\ii\infty}^{b+\ii\infty}
 -\frac{2\pi}{w}\cot\frac{\pi w}{2}
 \Gamma(w)\zeta(w)^2(t/A)^w\dd w.
\end{equation}

On the other hand, shift \eqref{eq:mellin} for $\eps=0$ to
$\Re s=-b$. The poles crossed give $\CC_0(\log(1/t))-t/4$.
The functional equation yields
\[
 \Gamma(s)\zeta(1+s)^2
 =-\frac{2\pi}{s}\cot\frac{\pi s}{2}
   \Gamma(-s)\zeta(-s)^2 A^s.
\]
Putting $w=-s$, including the orientation change, turns the remaining
integral precisely into \eqref{eq:median-MB}. This proves the exact
identity, including any possible flat real term. Equation
\eqref{eq:stokes-exact} is Theorem~\ref{thm:parity} integrated once.
\end{proof}

\begin{corollary}[Arithmetic Stokes data]
The positive Borel singularities lie at $An$, $n\ge1$. The logarithmic
jump accumulates the divisor counting function, and its Laplace image is
an absolutely convergent sum of exponential integrals. In particular,
\begin{equation}\label{eq:J-leading}
 J(t)=\frac{4\pi\ii}{A}\,t e^{-A/t}\bigl(1+O(t)\bigr)
 \qquad(t\to0^+).
\end{equation}
\end{corollary}
\begin{proof}
The local singularities and jump were determined explicitly. Integration
by parts in $E_1(A/t)$ gives its first term and an $O(t)$ relative
remainder. Every $n\ge2$ term is exponentially smaller; the bound
$d(n)\le n$ is already sufficient to sum that tail.
\end{proof}

\section{Uniform inverse crossover and the displaced core folds}
\label{sec:crossover}

\subsection{A convergent chart for the core, not for the full tail}
We use a scaled target $Y=X^2/2$ and set
\begin{equation}\label{eq:crossover-variables}
 u=X^{-1},\qquad \eps=\lambda u,\qquad w=u(L+\gamma).
\end{equation}
The scaled core equation is $G(w,\lambda,u)=0$, where
\begin{equation}\label{eq:G}
 G(w,\lambda,u)=
 \frac{\ua(\lambda u)e^{\lambda w}
       +\ua(-\lambda u)e^{-\lambda w}+\uz(\lambda u)}{\lambda^2}
 -\frac12.
\end{equation}
The value at $\lambda=0$ is removable. The cancellation $a_1=0$ is
what eliminates a term linear in $u$.

\begin{theorem}[Arcsinh crossover and its first corrections]
\label{thm:crossover}
Define the branches
\begin{equation}\label{eq:hs}
 h(\lambda)=\frac{2\arsinh(\lambda/2)}{\lambda},\quad h(0)=1,
 \qquad s(\lambda)=\sqrt{1+\lambda^2/4},\quad s(0)=1.
\end{equation}
On every compact subset of $|\lambda|<2$, the core inverse has a
holomorphic chart
\begin{equation}\label{eq:core-chart}
 L_c=\frac{h(\lambda)}u-\gamma+uB(u,\lambda),
\end{equation}
for sufficiently small $u$, with a genuinely convergent Taylor series
in $u$. Its first two coefficients are
\begin{equation}\label{eq:B-coeffs}
 B(0,\lambda)=-\frac{C_*+b\lambda^2}{s(\lambda)},
 \qquad \partial_uB(0,\lambda)=-2a_3\lambda^2.
\end{equation}
Equivalently,
\begin{equation}\label{eq:Lc-explicit}
 L_c=Xh(\lambda)-\gamma-
 \frac{C_*+b\lambda^2}{X\sqrt{1+\lambda^2/4}}
 -\frac{2a_3\lambda^2}{X^2}+O(X^{-3}).
\end{equation}
For real $\lambda$, the same assertion holds uniformly on any fixed
compact interval, using analytic neighborhoods of that interval.
\end{theorem}
\begin{proof}
At $u=0$,
\[
 G(w,\lambda,0)=\frac{\cosh(\lambda w)-1}{\lambda^2}-\frac12.
\]
The branch $w=h(\lambda)$ solves this equation because
$\cosh(2\arsinh(\lambda/2))=1+\lambda^2/2$.
Its $w$ derivative is
$\sinh(\lambda h)/\lambda=s(\lambda)$, which does not vanish
in $|\lambda|<2$ or at any real $\lambda$.
The holomorphic implicit function theorem therefore gives a unique
chart $w(u,\lambda)$ locally along this branch. Uniqueness glues the
local charts on compact sets. The absence of a linear $u$ term in
\eqref{eq:G} implies $w=h+u^2B$.

Writing $c_\lambda=1+\lambda^2/2$, the expansion is
\begin{align*}
 G(w,\lambda,u)
 &=\frac{\cosh(\lambda w)-1}{\lambda^2}-\frac12
   +u^2\{2b\cosh(\lambda w)+\delta_0\}\\
 &\quad+2a_3\lambda u^3\sinh(\lambda w)+O(u^4).
\end{align*}
Since $2b+\delta_0=C_*$, the $u^2$ coefficient gives
$sB(0,\lambda)+C_*+b\lambda^2=0$.
The $u^3$ coefficient gives
$s\partial_uB(0,\lambda)+2a_3\lambda\sinh(\lambda h)=0$.
Using $\sinh(\lambda h)=\lambda s$ proves
\eqref{eq:B-coeffs} and \eqref{eq:Lc-explicit}.
For a real compact interval $s(\lambda)>0$ throughout, so the same
argument supplies uniform neighborhoods and estimates.
\end{proof}

Every higher coefficient can be computed recursively from the analytic
equation \eqref{eq:G}. For example, if
$w=h+B_0u^2+B_1u^3+B_2u^4+\cdots$, then
\begin{equation}\label{eq:B2}
 B_2=-\frac{\tfrac12c_\lambda B_0^2
      +2b\lambda^2s(\lambda)B_0
      +\lambda^2(2a_4c_\lambda+z_4)}{s(\lambda)}.
\end{equation}
At $\lambda=0$ this reduces to $B_2=-C_*^2/2$, as required by
$L_c=\sqrt{X^2-2C_*}-\gamma$.

\subsection{Transport from the core to the actual positive inverse}
\begin{theorem}[First exponential correction in the crossover window]
\label{thm:actual-crossover}
Fix a compact real interval of $\lambda$ values. Let
$\eps=\lambda/X$, let $L_c$ be the core root of
Theorem~\ref{thm:crossover}, and put $t_c=e^{-L_c}$.
The unique positive solution of $\F_\eps(t)=X^2/2$ has
$L=\log(1/t)$ satisfying
\begin{equation}\label{eq:actual-crossover}
 L-L_c=\frac{\zeta(\eps)\zeta(-\eps)}
                 {\partial_L\CC_\eps(L_c)}t_c
       +O(t_c^2/X),
\end{equation}
uniformly as $X\to\infty$. Hence
\begin{equation}\label{eq:t-crossover}
 t=t_c-\frac{\zeta(\eps)\zeta(-\eps)}
                 {\partial_L\CC_\eps(L_c)}t_c^2
       +O(t_c^3/X).
\end{equation}
The denominator is $Xs(\lambda)+O(X^{-1})$.
\end{theorem}
\begin{proof}
Differentiate the scaled core equation with respect to $L$:
$\partial_L\CC_\eps(L_c)=XG_w(w_c,\lambda,u)$.
The chart gives $G_w=s(\lambda)+O(u^2)$, uniformly bounded away
from zero on the real compact set. Also $\partial_L^2\CC_\eps$
is uniformly bounded there. By Theorem~\ref{thm:collision},
\[
 \F_\eps(e^{-L})=\CC_\eps(L)+r_1(\eps)e^{-L}+O(e^{-2L}).
\]
The derivative in $L$ is positive and comparable to $X$ near $L_c$;
this follows either from the differentiated Mellin remainder or from
the positive defining sum. A residual estimate first gives
$L-L_c=O(t_c/X)$. Substitute this estimate back into the equation
and Taylor expand the core and $e^{-L}$ once. Then
\[
 0=\partial_L\CC_\eps(L_c)(L-L_c)+r_1(\eps)t_c+O(t_c^2).
\]
Use \eqref{eq:r1} and divide by the derivative. Exponentiation gives
\eqref{eq:t-crossover}; the square of $O(t_c/X)$ is absorbed by
the stated remainder. Positivity and uniqueness are supplied by
Proposition~\ref{prop:family}.
\end{proof}

A finite $1/X$ truncation of $L_c$ does \emph{not} isolate the
exponential correction: its omitted algebraic error is much larger.
Formula \eqref{eq:actual-crossover} uses the entire analytic core root,
not a short truncation of \eqref{eq:Lc-explicit}.

\subsection{Explicit critical-locus displacement for the complex core}
\begin{theorem}[Two moving simple folds of the normalized core]
\label{thm:folds}
Set
\begin{equation}\label{eq:Dfold}
 D=\delta_0-2b=3\gamma^2+6\gamma_1-\frac{\zeta(2)}2.
\end{equation}
Near $(w,\lambda,u)=(\pi/2,\pm2\ii,0)$ there are unique analytic
solutions of $G=G_w=0$. They satisfy
\begin{align}
 \lambda_\pm(u)&=\pm2\ii\{1+Du^2\}+O(u^4),
 \label{eq:fold-lambda}\\
 w_\pm(u)&=\frac\pi2\{1-Du^2\}+8a_3u^3+O(u^4).
 \label{eq:fold-w}
\end{align}
Each remains a simple fold: locally, the two core inverse branches have
a square-root interchange as $\lambda$ encircles the critical value.
This assertion concerns the singular core $\CC_\eps$, not the full
Lambert function at these complex scalings.
\end{theorem}
\begin{proof}
At the positive imaginary point,
\[
 G=G_w=0,\quad G_{ww}=-1,\quad G_\lambda=\ii/2,
 \quad G_{w\lambda}=\ii\pi/4.
\]
The Jacobian of $(G,G_w)$ with respect to $(\lambda,w)$ is
nonzero. The implicit function theorem gives the critical locus.
The $u^2$ term of $G$ at the point is $\delta_0-2b=D$, so
$(\ii/2)(\lambda_+(u)-2\ii)+Du^2=O(u^3)$.
This gives the $u^2$ coefficient of \eqref{eq:fold-lambda}.
The $u^2$ coefficient in $G_w=0$ gives
$w_+(u)-\pi/2=-\pi Du^2/2+O(u^3)$.

The order-$u^3$ term in $G$ is proportional to
$\sinh(\lambda w)$ and vanishes at the limiting fold; hence
$\lambda_+$ has no $u^3$ term. Its $w$ derivative at that point
is $2a_3\lambda^2\cosh(\lambda w)=8a_3$, giving the displayed
$8a_3u^3$ term in $w_+$. Evenness of $G$ in $\lambda$ gives the
negative branch. Finally $G_{ww}\ne0$ and $G_\lambda\ne0$
persist. Taylor expansion and the Weierstrass preparation theorem
reduce the equation to a nondegenerate quadratic in the local $w$
coordinate; its discriminant has a simple zero in $\lambda$.
That is precisely the asserted square-root interchange.
\end{proof}

The fold location $\lambda=\pm2\ii$ is distinct from the secondary
resonances $2\pi\ii k$ in the single-polylogarithm crossover discussed
in the earlier draft \cite{prior-folds}. The distinction follows from
the different leading functions: $\cosh(\lambda w)-1$ here, rather
than $(1-e^{-\lambda})/\lambda$ there.

\section{Exact inverse-sector coefficients: signs, Catalan degrees, and divergence}
\label{sec:coefficients}
At exact resonance it is preferable to absorb the constant $C_*$ into
the large variable. In this section only, set
\begin{equation}\label{eq:exact-core-vars}
 X=\sqrt{2(Y-C_*)},\qquad u=X^{-1},\qquad
 q=e^{\gamma-X},\qquad L=X-\gamma+V.
\end{equation}
This $X$ differs from the scaled-target variable of
Section~\ref{sec:crossover}; both definitions are explicitly tied to
their targets. The positive branch is used, and $Y$ is sufficiently
large. Then $t=qe^{-V}$. Define the formal series including the linear
term
\begin{equation}\label{eq:R-all}
 R(q)=\sum_{n\ge1}r_nq^n,\qquad r_n=r_n(0).
\end{equation}
The formal inverse equation is
\begin{equation}\label{eq:formal-inverse}
 V+\frac u2V^2+uR(qe^{-V})=0.
\end{equation}
Here $q$ and $u$ are independent formal variables for coefficient
extraction. They are specialized to the physical relation
\eqref{eq:exact-core-vars} only afterward.

\begin{theorem}[Finite coefficient formula and sharp polynomial structure]
\label{thm:inverse-coefficients}
There is a unique solution
\begin{equation}\label{eq:V-series}
 V(q,u)=\sum_{n\ge1}C_n(u)q^n\in q\Q[u][[q]]
\end{equation}
of \eqref{eq:formal-inverse}. For every $n\ge1$ it is given by
\begin{equation}\label{eq:C-formula}
 C_n(u)=-\sum_{k=1}^{n}\frac{u^k}{k!}[q^n]R(q)^k
 \sum_{j=0}^{k-1}\binom{k-1}{j}n^{k-1-j}(k)_j\left(\frac u2\right)^j.
\end{equation}
Its structural properties are
\begin{align}
 \deg C_n&=2n-1,&
 [u^{2n-1}]C_n&=\frac{(-1)^{n+1}\Cat_{n-1}}{2^{3n-1}},
 \label{eq:C-degree}\\
 (-1)^{n+1}C_n(u)&\in\Q_{\ge0}[u]\setminus\{0\},
 \label{eq:C-sign}\\
 \operatorname{ord}_{u=0}C_n
 &=\begin{cases}1,&n=1\text{ or }n\text{ even},\\2,&n\ge3\text{ odd},\end{cases}
 \label{eq:C-valuation}\\
 [u]C_n&=-r_n,&
 [u^2]C_n&=\frac n4r_{n-1}\quad(n\ge3\text{ odd}).
 \label{eq:C-lowest}
\end{align}
In \eqref{eq:C-sign} every nonzero coefficient is strictly positive
after the indicated sign change. In particular every inverse sector is
present, although every odd forward coefficient beyond the linear term
vanishes.
\end{theorem}
\begin{proof}
The constant term in $q$ must be $V=0$; the derivative with respect to
$V$ at $q=V=0$ is $1$. Equating coefficients recursively gives a unique
solution over $\Q[u]$, since each $r_n$ is rational.

For an explicit formula introduce an auxiliary variable $\eta$ and
write
\[
 V=-\eta u\frac{R(qe^{-V})}{1+uV/2}.
\]
Formal Lagrange inversion in $\eta$ yields
\begin{equation}\label{eq:eta-Lagrange}
 V=\sum_{k\ge1}\frac{(-\eta u)^k}{k!}
 \left.\left(\frac{\partial}{\partial v}\right)^{k-1}
 \left\{R(qe^{-v})^k(1+uv/2)^{-k}\right\}\right|_{v=0}.
\end{equation}
Because $R$ has positive $q$ valuation, only $k\le n$ contributes
to $q^n$. Thus setting $\eta=1$ is legitimate coefficientwise, even
though no analytic convergence of $R$ is being assumed.
For the $q^n$ coefficient,
$[q^n]R(qe^{-v})^k=e^{-nv}[q^n]R(q)^k$.
Leibniz' rule gives
\[
 \left.\partial_v^{k-1}\{e^{-nv}(1+uv/2)^{-k}\}\right|_0
 =(-1)^{k-1}\sum_{j=0}^{k-1}
 \binom{k-1}{j}n^{k-1-j}(k)_j(u/2)^j.
\]
Combining signs proves \eqref{eq:C-formula}.

Since $k+j\le2k-1\le2n-1$, the degree is at most $2n-1$.
The only term of that degree has $k=n$, $j=n-1$, and uses
$[q^n]R^n=r_1^n=(-1/4)^n$. Hence its coefficient is
\[
 -\frac{(-1/4)^n}{n!}\frac{(n)_{n-1}}{2^{n-1}}
 =\frac{(-1)^{n+1}}{2^{3n-1}}
   \frac{(2n-2)!}{n!(n-1)!},
\]
which is exactly \eqref{eq:C-degree}, and is nonzero.

For the sign assertion, every nonzero monomial contributing to
$[q^n]R^k$ consists of factors $r_1<0$ and $r_{2m}>0$.
The number of $r_1$ factors has the same parity as $n$.
Consequently $(-1)^n[q^n]R^k\ge0$. Every coefficient of the inner
polynomial in \eqref{eq:C-formula} is nonnegative, so no cancellation
can reverse this sign. The degree computation ensures nonvanishing.
The coefficient of $u$ comes only from $k=1$, and is $-r_n$.
For odd $n\ge3$ this vanishes. The coefficient of $u^2$ comes only
from $k=2,j=0$ and equals $-n[q^n]R^2/2$.
For odd $n\ge3$, $[q^n]R^2=2r_1r_{n-1}=-r_{n-1}/2$.
This proves \eqref{eq:C-lowest} and the exact valuations.
\end{proof}

The first three coefficients, useful as normalization checks, are
\begin{align}
 C_1(u)&=\frac u4,\label{eq:C1}\\
 C_2(u)&=-\frac u{288}-\frac{u^2}{16}-\frac{u^3}{32},
 \label{eq:C2}\\
 C_3(u)&=\frac{u^2}{384}+\frac{7u^3}{288}
         +\frac{3u^4}{128}+\frac{u^5}{128}.
 \label{eq:C3}
\end{align}
The supplied exact verifier independently substitutes all coefficients
through $q^8$ into \eqref{eq:formal-inverse}. The proof, rather than
that finite computation, establishes the all-order assertions.

\subsection{Exact Borel radius of the inverse-sector series}
\begin{theorem}[Gevrey order and sharp Borel radius]
\label{thm:inverse-gevrey}
For every $U>0$, the polynomials in
Theorem~\ref{thm:inverse-coefficients} satisfy, uniformly for
$|u|\le U$,
\begin{equation}\label{eq:gevrey-bound}
 |C_n(u)|\le K_U A^{-n}\Gamma(n)e^{K_U\sqrt n}.
\end{equation}
For every fixed real $u>0$, the ordinary series
$\sum C_n(u)q^n$ has radius zero, whereas its shifted Borel transform
\begin{equation}\label{eq:inverse-borel}
 \widehat V_u(\xi)=\sum_{n\ge1}\frac{C_n(u)}{\Gamma(n)}\xi^{n-1}
\end{equation}
has radius exactly $A=4\pi^2$. In particular it is a Gevrey-one
series, but not a convergent exponential-sector chart at fixed $u>0$.
The point $\xi=-A$ is a singularity of the Borel germ.
\end{theorem}
\begin{proof}
Formula \eqref{eq:even-coefficients} gives a constant $K>0$ with
$|r_m|\le K A^{-m}(m-1)!$ for all $m\ge1$.
For a composition $m_1+\cdots+m_k=n$ with $m_i\ge1$,
$\prod_i(m_i-1)!\le(n-k)!$. Counting the compositions gives
\begin{equation}\label{eq:composition-bound}
 |[q^n]R(q)^k|
 \le K^k A^{-n}(n-k)!\binom{n-1}{k-1}
 =K^k A^{-n}\frac{(n-1)!}{(k-1)!}.
\end{equation}
Also $(k)_j\le(2n)^j$ for $0\le j\le k-1$, so the absolute
value of the inner sum in \eqref{eq:C-formula} is at most
$n^{k-1}(1+U)^{k-1}$. Therefore
\[
 |C_n(u)|\le A^{-n}\Gamma(n)KU
 \sum_{j=0}^{n-1}\frac{\{KUn(1+U)\}^j}{(j+1)!j!}.
\]
The elementary bound
$\sum_{j\ge0}x^j/(j!)^2\le e^{2\sqrt x}$ for $x\ge0$
follows by comparison with the product
$e^{\sqrt x}e^{\sqrt x}$. This proves \eqref{eq:gevrey-bound}.
For any $A'<A$, its right-hand side is at most
$K_{U,A'}(A')^{-n}\Gamma(n)$, which is the usual Gevrey-one
estimate. In particular the Borel radius is at least $A$.

Fix $u>0$. For even $n$, all terms in $C_n(u)$ have the same
negative sign, and the $k=1$ term gives
\[
 |C_n(u)|\ge u r_n
 =\frac{4u\Gamma(n)}{nA^n}\zeta(n)^2.
\]
This lower bound proves radius zero for the ordinary series and Borel
radius at most $A$. Thus the latter radius is exactly $A$.
Finally \eqref{eq:C-sign} shows that $\widehat V_u(-\xi)$ has
nonnegative Taylor coefficients. A power series with nonnegative
coefficients and finite radius $A$ cannot be analytic at its positive
endpoint $A$: otherwise its Taylor expansion about a point just below
$A$, with nonnegative derivatives, can be rearranged to make the
original series converge at a point beyond $A$. This contradiction
is the elementary positive-coefficient form of Pringsheim's argument.
Applied to $\widehat V_u(-\xi)$, it proves singularity at $-A$.
\end{proof}

\begin{remark}[What the sharp radius does not prove]
The theorem locates the nearest guaranteed negative-axis singularity and
fixes the radius of the Borel germ. It does not, by itself, prove global
analytic continuation of the inverse Borel transform, classify all its
singularities, or identify every alien derivative. The lateral inverse
functions in Section~\ref{sec:inverse-stokes} are defined analytically
from the \emph{exact forward lateral functions}, so their jump does not
depend on those stronger unproved assertions.
\end{remark}

\subsection{The actual inverse and finite-sector error control}
\begin{theorem}[Physical specialization of the formal coefficients]
\label{thm:physical-sectors}
Let $t=g_0(Y)$, with the physical variables \eqref{eq:exact-core-vars},
and let $V=\log(q/t)$. For every fixed $N\ge1$,
\begin{equation}\label{eq:physical-V}
 V=\sum_{n=1}^{N}C_n(1/X)q^n+O(q^{N+1}/X),
 \qquad X\to+\infty.
\end{equation}
In particular,
\begin{equation}\label{eq:t-first}
 t=q\left(1-\frac{q}{4X}+O(q^2/X)\right),
 \qquad
 \frac A t=\frac A q+\frac{A}{4X}+O(q/X).
\end{equation}
\end{theorem}
\begin{proof}
The exact expansion \eqref{eq:asymptotic} at $\eps=0$, with a
fixed finite remainder order, gives
\[
 0=XV+\frac12V^2+\sum_{m=1}^{N}r_mq^m e^{-mV}+O(q^{N+1})
\]
near $V=0$. Monotonicity first gives $V=O(q/X)$, for example by
comparing the residual at $V=0$ with the derivative, which is
$X+O(q)$. Substitute
$V_N=\sum_{n\le N}C_n(1/X)q^n$.
The formal coefficient identity cancels every power through $q^N$.
Since $N$ is fixed and the coefficients are polynomials bounded as
$1/X\to0$, the remaining residual is $O(q^{N+1})$.
The derivative of the exact equation is bounded below by $X/2$ in
this neighborhood, giving the stated $O(q^{N+1}/X)$ correction.
Using $C_1=1/(4X)$ and exponentiating proves \eqref{eq:t-first}.
\end{proof}

This is an asymptotic statement at every fixed sector order. It is not
an assertion that the infinite divergent series may be evaluated at the
physical pair $(q,1/X)$ by ordinary summation.

\section{Exact inverse Stokes transport and the nested exponential scale}
\label{sec:inverse-stokes}
The forward jump in \eqref{eq:stokes-exact} is explicit. We now transport
it through functional inversion without silently linearizing away the
higher nonlinear sectors.

\subsection{A convergent Lagrange formula in the Stokes parameter}
Let $g_0$ denote the positive inverse of $\F_0$. For $Y$ sufficiently
large, let $g_\eta(Y)$ be the local inverse near $t=g_0(Y)$ of
\begin{equation}\label{eq:Feta}
 \F_\eta(z)=\F_0(z)+\frac\eta2J(z).
\end{equation}
The analytic continuations used here lie in a small disk around positive
$t$; $E_1(An/z)$ uses the continuation from positive $z$.
For $\eta=\pm1$ this agrees with the two forward lateral functions.

\begin{theorem}[All-order analytic transport of the arithmetic jump]
\label{thm:inverse-transport}
For sufficiently large real $Y$, the inverse $g_\eta(Y)$ is holomorphic
in $\eta$ on a disk of radius greater than $1$, and
\begin{equation}\label{eq:inverse-transport}
 g_\eta(Y)=g_0(Y)+\sum_{k\ge1}\frac{(-\eta)^k}{k!}
 \partial_Y^{k-1}\left[
 g_0'(Y)\left\{\frac{J(g_0(Y))}{2}\right\}^{k}\right].
\end{equation}
The series is actually convergent at $\eta=\pm1$. Consequently
\begin{align}
 g_+(Y)-g_-(Y)
 &=-2\sum_{j\ge0}\frac1{(2j+1)!}
 \partial_Y^{2j}\left[
 g_0'(Y)\left\{\frac{J(g_0(Y))}{2}\right\}^{2j+1}\right],
 \label{eq:inverse-jump-all}\\
 \frac{g_+(Y)+g_-(Y)}2-g_0(Y)
 &=\sum_{j\ge1}\frac1{(2j)!}
 \partial_Y^{2j-1}\left[
 g_0'(Y)\left\{\frac{J(g_0(Y))}{2}\right\}^{2j}\right].
 \label{eq:inverse-average-all}
\end{align}
These formulas retain the entire arithmetic divisor sum in $J$, not
only its first action.
\end{theorem}
\begin{proof}
First establish a domain in which the series is more than formal.
Let $t=g_0(Y)$ and $W=\log(1/t)+\gamma$. On a disk
$|z-t|\le c t^2$, with fixed sufficiently small $c>0$, the differentiated
Mellin expansion gives
\[
 \F_0'(z)=-\frac{\log(1/z)+\gamma}{z}-\frac14+O(z),
 \qquad \F_0''(z)=O(W/t^2).
\]
The disk stays in the right half-plane. Taylor's theorem then shows
that on its boundary
$|\F_0(z)-Y|\ge c_1 Wt$, and that $\F_0(z)-Y$ has one simple
zero inside. The exponential integral representation gives uniformly
\[
 |J(z)|\le c_2 t e^{-A/t},
\]
since $1/z=1/t+O(1)$ throughout the disk; the arithmetic tail is
summable uniformly there.
Put $\rho(t)=\exp(3A/(4t))$. Then
$\rho(t)|J(z)|=O(t e^{-A/(4t)})=o(Wt)$.
Rouch\'e's theorem shows that for every $|\eta|\le\rho(t)$ the
perturbed equation has exactly one zero in the disk, counted with
multiplicity. It is therefore simple, and the implicit function theorem
makes it holomorphic in $\eta$. In particular its Taylor series
converges at $\eta=\pm1$ for small enough $t$.

To compute that Taylor series, use the local coordinate
$v=\F_0(z)$ and write $z=g_0(v)$, extending the inverse locally in
$Y$. Equation \eqref{eq:Feta} becomes
\[
 v=Y-\eta h(v),\qquad h(v)=J(g_0(v))/2.
\]
Lagrange--B\"urmann applied to the analytic function $g_0(v)$ gives
\eqref{eq:inverse-transport}. Taking its odd and even parts in
$\eta$ gives the other two identities. This argument uses ordinary
local analytic inversion in an auxiliary parameter, not an assumed
convergence of the factorial $q$-sector series.
\end{proof}

\subsection{The first two inverse Stokes effects}
\begin{theorem}[Nested jump and exact leading averaging defect]
\label{thm:nested}
For $t=g_0(Y)\to0^+$, $W=\log(1/t)+\gamma$, and $A=4\pi^2$,
\begin{align}
 g_+(Y)-g_-(Y)
 &=\frac{4\pi\ii}{A}\frac{t^2}{W}e^{-A/t}\{1+O(t)\},
 \label{eq:nested-jump}\\
 \frac{g_+(Y)+g_-(Y)}2-g_0(Y)
 &=-\frac{t^2}{W^2}e^{-2A/t}\{1+O(t)\}.
 \label{eq:nested-bias}
\end{align}
In the physical variables of \eqref{eq:exact-core-vars},
\begin{align}
 g_+(Y)-g_-(Y)
 &=\frac{\ii}{\pi}\frac{q^2}{X}
   \exp\left(-\frac A q-\frac{A}{4X}\right)\{1+O(q)\},
 \label{eq:nested-q}\\
 \frac{g_+(Y)+g_-(Y)}2-g_0(Y)
 &=-\frac{q^2}{X^2}
   \exp\left(-\frac{2A}{q}-\frac{A}{2X}\right)\{1+O(q)\}.
 \label{eq:bias-q}
\end{align}
Thus inversion exposes a doubly exponential Stokes level, with leading
action $A/q=Ae^{X-\gamma}$.
\end{theorem}
\begin{proof}
The exact forward expansion and \eqref{eq:J-leading} imply
\begin{equation}\label{eq:derivatives-leading}
 g_0'(Y)=-\frac{t}{W}\{1+O(t/W)\},\qquad
 J(t)=\frac{4\pi\ii}{A}t e^{-A/t}\{1+O(t)\}.
\end{equation}
The first term of \eqref{eq:inverse-jump-all} is
$-g_0'(Y)J(t)$, giving \eqref{eq:nested-jump}.
For rigor in discarding the remaining terms, the proof of
Theorem~\ref{thm:inverse-transport} bounds $g_\eta(Y)-t$ by
$ct^2$ on $|\eta|=\rho(t)$. Its Taylor coefficients are therefore
bounded by $ct^2\rho(t)^{-k}$. The remaining odd terms begin at
$k=3$, and are $O(t^2e^{-9A/(4t)})$, which is negligible relative
to the first term.

The first term of \eqref{eq:inverse-average-all} is
\begin{equation}\label{eq:bias-leading-exact}
 \frac18\partial_Y\{g_0'(Y)J(t)^2\}.
\end{equation}
The remaining even terms begin at $k=4$ and are
$O(t^2e^{-3A/t})$, again negligible. Since
$\partial_Y=g_0'(Y)\partial_t$, the leading product inside
\eqref{eq:bias-leading-exact} is
\[
 \frac{16\pi^2}{A^2}\frac{t^3}{W}e^{-2A/t}.
\]
Differentiating its exponential and multiplying by $-t/W$ gives
\[
 -\frac{32\pi^2}{A}\frac{t^2}{W^2}e^{-2A/t}.
\]
Division by $8$ leaves $-4\pi^2/A=-1$. Derivatives of the
algebraic prefactor, and the differentiated remainder terms, produce a
relative $O(t)$ error. These derivative estimates follow either by
integration by parts in $E_1$ with differentiated remainders or by its
exact derivative; they are not obtained by differentiating an arbitrary
uncontrolled $O$ term. This proves \eqref{eq:nested-bias}.

Finally Theorem~\ref{thm:physical-sectors} gives
$t=q\{1-q/(4X)+O(q^2/X)\}$, $W=X+O(q/X)$, and
$A/t=A/q+A/(4X)+O(q/X)$. Substitute these expressions in the
two preceding formulas. The exponent error is $O(q/X)$, so it can
be exponentiated into a relative $1+O(q)$ factor, proving
\eqref{eq:nested-q}--\eqref{eq:bias-q}.
\end{proof}

\begin{remark}[A normalization safeguard]
Replacing $t$ merely by $q$ \emph{inside} $e^{-A/t}$ loses the
specific algebraic exponential multiplier $e^{-A/(4X)}$.
Although that multiplier tends to $1$, it matters when resolving
successive slow corrections within the same nested action.
At the same time, arithmetic averaging of $g_+$ and $g_-$ is not the
same operation as inverting the arithmetic average of $\F_+$ and
$\F_-$. Equation \eqref{eq:nested-bias} quantifies that difference.
This does not contradict multiplicativity properties of a properly
defined nonlinear median summation operator: no such operator is being
identified with a plain arithmetic average of the inverse branches.
\end{remark}

\section{Verification and reproducibility}
\label{sec:verification}
The proofs above are independent of the verification program. The
program is intended to catch normalization, sign, indexing, branch,
and transcription errors before integration.

\subsection{Exact checks}
The SymPy portion constructs \eqref{eq:C-formula} through $n=8$ using
rational zeta values. It then substitutes the resulting polynomials
into \eqref{eq:formal-inverse} by an independent recurrence for the
coefficients of $e^{-mV}$:
\[
 E_0=1,\qquad
 E_n=-\frac mn\sum_{j=1}^n jC_j E_{n-j}.
\]
Every residual coefficient through $q^8$ is checked to be exactly zero.
The program also checks the polynomial degrees, valuations, all
coefficient signs, the Catalan leading coefficients, and the
odd-sector $u^2$ coefficients. A separate symbolic multiplication
checks $a_0,a_1,a_2,a_3$ in the centered normalizer. These are finite
exact checks, not formal proofs of the infinite family.

\subsection{Numerical diagnostics}
The mpmath portion compares the two sides of the reflected coefficient
identity at real and complex orders; checks the second-jet formula at
three odd indices; verifies pole coalescence; solves the scaled core
equation and its two-variable critical equations; and compares the
forward asymptotic expansion with the absolutely convergent
$\Li_1$ sum. An optional SciPy calculation independently compares
\eqref{eq:median-MB} with the original Lambert sum at $t=1$ and
$t=3$. It truncates the vertical integral at height $55$ and records
the adaptive quadrature error estimate. That estimate is not an
interval certificate.

For the most delicate inverse checks, the program solves
$\F_0(z)+J(z)/2=\F_0(t)$ at $t=1,1/2,1/4$ with increased precision.
Newton iteration uses the independently differentiated convergent
series
\begin{equation}\label{eq:numeric-derivatives}
 \F_0'(t)=-\sum_{a\ge1}\frac{e^{-at}}{1-e^{-at}},\qquad
 J'(t)=\frac{4\pi\ii}{t}\sum_{n\ge1}d(n)e^{-An/t}.
\end{equation}
The conjugate root is the lower lateral inverse. The recorded ratios
test both the first-action imaginary splitting and the much smaller
second-action real bias. Near this nested scale, a default generic
secant solver is not a reliable residual certificate; explicit Newton
steps and a final residual check are used instead.

\begin{center}
\begin{tabular}{@{}rrrr@{}}
\toprule
$t$ & Working digits & Jump / leading term & Bias / leading term\\
\midrule
$1$   & $80$  & $0.686714522854$ & $0.494775281614$\\
$1/2$ & $125$ & $0.900289105501$ & $0.829028586173$\\
$1/4$ & $205$ & $0.963296642480$ & $0.938079390839$\\
\bottomrule
\end{tabular}
\end{center}
These are ratios to the asymptotic leading terms in
\eqref{eq:nested-jump} and \eqref{eq:nested-bias}, not relative errors
in the exact identities. Their approach to $1$ is a diagnostic for the
stated small-$t$ laws. A recorded zero root residual at finite working
precision means only that the residual rounded to zero; it is not an
exact certificate.

\subsection{How to reproduce}
From the package directory, run
\begin{verbatim}
python verification/verify.py --output verification/results.json
pdflatex -interaction=nonstopmode -halt-on-error article.tex
pdflatex -interaction=nonstopmode -halt-on-error article.tex
\end{verbatim}
The file \texttt{verification/results.json} records the actual run,
including software versions and the scope of each check. The flag
\texttt{--skip-lateral} omits only the higher-precision lateral-root
diagnostics. SciPy is optional; without it, the independent Mellin
quadrature is explicitly marked as not run. No network access is
needed by the verification program or the TeX build.

\section{Source audit, corrections, and integration boundaries}
\label{sec:audit}

\subsection{What was inspected}
The repository was anchored at commit
\texttt{6207c3d40047fd8d1400a4d509fee71c9a3a8132}, whose recorded
message was ``New research reports.'' The inspection included the
Polylogarithms overview, the Transseries overview and its
series-and-transseries overview, the incoming directory inventory and
README, the manuscript chapter inventory, and the normalized
positive--negative derivative-bridge chapter. Three library research
drafts were used for an abstract-level comparison and selected
statement checks. The ZIP byte contents were not downloaded; no
byte-identity claim between those library drafts and individual ZIP
members is made. This is a targeted research audit, not a validation
of the entire repository.

\subsection{No unsupported source-specific correction}
No new erroneous theorem in the inspected repository material has been
established by this work. In particular, the convergence claims in the
selected prior drafts concern holomorphic analytic tails and do not
conflict with Theorem~\ref{thm:inverse-gevrey}. We therefore provide no
patch that weakens or rewrites those existing theorems.

Two \emph{proposed integration safeguards} are mathematically substantive.
First, replacing a single polylogarithm by the Lambert superposition
\eqref{eq:family} destroys the analyticity of the ordinary tail at the
origin: \eqref{eq:even-coefficients} proves radius zero. Any extension
that asserted an ordinary convergent inverse-sector series for this
new family at fixed $u>0$ would be false, and
Theorem~\ref{thm:inverse-gevrey} supplies its correct Gevrey replacement.
Second, interpreting the arithmetic mean of the lateral inverse
branches as the real inverse would be false, with the explicit
nonzero counterterm \eqref{eq:nested-bias}. Neither warning is
attributed to an inspected source that did not make the offending claim.

\subsection{Recommended integration shape}
The contribution is designed as a standalone research package, with
cross-links from Polylogarithms and Transseries. A suitable topic label is
\texttt{resonant-lambert-inverse-transseries}. The Borel and
Mellin statements should cite the established Lambert literature rather
than replace it. The new package should retain separate status labels
for conventional proofs, exact finite checks, ordinary numerical
diagnostics, and future proof-assistant work. No theorem here has been
compiled in Lean, and no numerical output is presented as a rigorous
interval enclosure.

The \texttt{AUDIT.md}, \texttt{PROVENANCE.md}, and \texttt{CLAIMS.md}
files summarize these distinctions outside the TeX article. The package
does not alter the canonical manuscript or publish any repository commit.

\section{Further research questions}
\label{sec:questions}
The following problems are concrete continuations, not premises used
in the preceding proofs.

\subsection{Sharp large-order constants of the inverse sectors}
Theorem~\ref{thm:inverse-gevrey} fixes the exponential Borel scale but
not the full large-order coefficient law. A natural conjecture, for
fixed $u>0$, is
\begin{align}
 C_{2m}(u)&\sim-\frac{4u\Gamma(2m)}{2mA^{2m}}
                         \cosh(Au/4),\label{eq:conj-even}\\
 C_{2m+1}(u)&\sim\frac{4u\Gamma(2m+1)}{(2m+1)A^{2m+1}}
                         \sinh(Au/4).\label{eq:conj-odd}
\end{align}
These are \emph{unproved in this article}. For each fixed number of
factors in \eqref{eq:C-formula}, the contribution of one large even
index and the remaining linear indices leads to the displayed
hyperbolic series. A proof must control the sum over the number of
factors uniformly with $m$; the current $e^{K\sqrt n}$ upper bound
is not sufficient for that interchange. The conjecture would distinguish
the amplitudes at $+A$ and $-A$ and sharpen the reversion calculus.

\subsection{Full inverse Borel continuation}
Does $\widehat V_u$ continue along every path avoiding an explicit
additive arithmetic action set? Can one describe its positive-axis
singularity directly from \eqref{eq:C-formula}, rather than from
analytic inversion of the forward lateral functions? The answer would
turn the present exact transport identities into a more complete
resurgent classification of the formal inverse itself.

\subsection{Uniform Stokes transport away from exact resonance}
The hypergeometric kernel \eqref{eq:general-borel} is jointly analytic
in the order parameter. Deriving its full discontinuity in closed
hypergeometric coordinates, and then resolving the scaled regime
$\eps=\lambda/X$, should connect the arcsinh crossover to moving
arithmetic Stokes amplitudes. All branch choices and the distinction
between core and full function must remain explicit.

\subsection{From core folds to full complex Lambert folds}
Theorem~\ref{thm:folds} concerns only the singular core. One may ask
whether suitable sectors permit a full-function deformation of those
critical loci, and whether the resulting inverse is uniformly described
by a quadratic local normal form with an exponentially small
arithmetic displacement. Direct substitution into a complex scaled
path is not enough: the full Lambert sum must remain in a domain of
analytic continuation with controlled remainders.

\subsection{Character pairs and higher collision multiplicity}
The two-character family in \cite{BD2025} permits the same inverse
questions with twisted divisor weights. Primitive character parities
should determine the vanishing order of perturbative sectors. Another
extension is the product-lattice sum over $r$ positive integers, whose
Mellin kernel contains $r$ zeta factors. At equal indices it has an
$(r+1)$-fold pole and an $r$-th degree logarithmic core. Its inverse
would naturally introduce $e^{-cY^{1/r}}$ sectors and nested Stokes
levels. Explicit centered order-jet formulas are a tractable first step.

\subsection{Certified computation and formal proof organization}
An interval implementation of the logarithmic Borel integral would
provide a practical check of exponentially small corrections without
subtraction of enormous nearly equal terms. On the formal side, the
coefficient formula, sign theorem, Catalan degree, and finite
composition estimates are natural initial targets. Mellin contour
shifts, parameter-uniform asymptotics, and analytic Stokes transport
form a separate, substantially larger analytic formalization task.

\section*{Conclusion}
\addcontentsline{toc}{section}{Conclusion}
The double Lambert specialization connects classical polylogarithms,
Gamma--Stieltjes order calculus, arithmetic Borel singularities, and
nonlinear transseries inversion in a single exactly normalized example.
The triple pole creates a quadratic logarithmic core; its scaled inverse
has an arcsinh crossover. Trivial-zero parity produces a full missing
forward sector, restored both by the second order jet and, through a
different mechanism, by nonlinear reversion. The reversion coefficients
have an exact rational formula, Catalan top degree, alternating signs,
and a sharp Borel scale despite ordinary divergence. Finally, the exact
arithmetic jump survives inversion as a nested exponential hierarchy,
and the first real averaging defect has the particularly simple
coefficient $-1$. These results provide a focused bridge between the
two repository areas while keeping classical input, proposed refinements,
proofs, tests, and unresolved extensions separate.

\begin{thebibliography}{20}
\bibitem{BD2025}
D. Broadhurst and D. Dorigoni,
\emph{Resurgent Lambert series with characters},
arXiv:2507.21352, version 1 (28 July 2025), especially Sections 2.3 and 3.
\url{https://arxiv.org/html/2507.21352v1}.

\bibitem{BD2026}
D. Broadhurst and D. Dorigoni,
\emph{Resurgent Lambert series from Feynman and beyond},
arXiv:2607.14020v1 (15 July 2026), Section 4, equations (30)--(31).
\url{https://arxiv.org/html/2607.14020v1}.

\bibitem{DK2021}
D. Dorigoni and A. Kleinschmidt,
\emph{Resurgent expansion of Lambert series and iterated Eisenstein integrals},
Communications in Number Theory and Physics \textbf{15} (2021), 1--57;
arXiv:2001.11035.
\url{https://arxiv.org/abs/2001.11035}.

\bibitem{DLMFpolylog}
NIST Digital Library of Mathematical Functions, Section 25.12,
\emph{Polylogarithms}, especially equation 25.12.10.
\url{https://dlmf.nist.gov/25.12}.

\bibitem{DLMFzeta}
NIST Digital Library of Mathematical Functions, Section 25.4,
\emph{Reflection Formulas} for the Riemann zeta function.
\url{https://dlmf.nist.gov/25.4}.

\bibitem{DLMFgamma}
NIST Digital Library of Mathematical Functions, Section 5.7,
\emph{Series Expansions} for the Gamma function.
\url{https://dlmf.nist.gov/5.7}.

\bibitem{repo-polylog}
V. Reshetnikov, \emph{ProveIt: Analysis/Polylogarithms}, repository
overview, inspected 12 October 2026.
\url{https://github.com/VladimirReshetnikov/ProveIt/tree/6207c3d40047fd8d1400a4d509fee71c9a3a8132/Analysis/Polylogarithms}.

\bibitem{repo-bridge}
V. Reshetnikov, \emph{ProveIt}, manuscript chapter
\texttt{08-bridges.tex}, ``One normalized functional equation for all
derivative bridges,'' inspected at the same commit.
\url{https://github.com/VladimirReshetnikov/ProveIt/blob/6207c3d40047fd8d1400a4d509fee71c9a3a8132/Analysis/Polylogarithms/docs/manuscript/chapters/08-bridges.tex}.

\bibitem{repo-transseries}
V. Reshetnikov, \emph{ProveIt: Analysis/Transseries}, overview and
\texttt{docs/series-and-transseries/README.md}, inspected 12 October 2026.
\url{https://github.com/VladimirReshetnikov/ProveIt/tree/6207c3d40047fd8d1400a4d509fee71c9a3a8132/Analysis/Transseries}.

\bibitem{repo-incoming}
V. Reshetnikov, \emph{ProveIt: docs/incoming}, inventory and README,
inspected at the same commit.
\url{https://github.com/VladimirReshetnikov/ProveIt/tree/6207c3d40047fd8d1400a4d509fee71c9a3a8132/docs/incoming}.

\bibitem{prior-jets}
\emph{Resonant Polylogarithmic Jets and Convergent Inverse Transseries:
Gamma normal forms, optimal dilation filters, arithmetic sector selection,
and exact order tuning}, research draft prepared for V. Reshetnikov,
11 October 2026. Library source \texttt{resonant\_polylog\_transseries.tex};
abstract and selected normal-form statements inspected.

\bibitem{prior-folds}
\emph{Resonant Polylogarithm Inversion: Uniform Order Crossover,
Gamma--Zeta Fold Loci, and Convergent Exponential Sectors}, research draft
prepared for V. Reshetnikov, 11 October 2026. Library source
\texttt{article(20261012-045040).tex}; abstract and selected core statements
inspected.

\bibitem{prior-fermi}
\emph{Resonant Polylogarithm Transseries: Gamma Borel Kernels,
Cyclotomic Order Jets, and Sharp Dilogarithm Reversion}, research draft
prepared for V. Reshetnikov, 11 October 2026. Library source
\texttt{article(20261012-044824).tex}; abstract and selected convergence
statements inspected.
\end{thebibliography}
\end{document}
